\documentclass[12pt]{article}
\usepackage[utf8]{inputenc}
\usepackage[T1]{fontenc}
\usepackage{hyperref}
\title{Atong Language Resources for K-12 Curriculum Development}
\author{Compiled by Claude Sonnet 5.5 (claude-sonnet-5-5) for Andr\'as Kornai}
\date{}
\begin{document}
\maketitle

\section{Introduction}
Atong (ISO 639-3: aot) is a Sino-Tibetan (Bodo-Garo) language spoken in Meghalaya and Assam in northeast India and in neighbouring Bangladesh, and is closely related to Garo. It is written in the Latin script and is classed as endangered. Only a small number of scholarly resources turned up in search; no school primers, textbooks or community-produced teaching materials were found.

\section{Primary Educational Resources}
No primers, textbooks or school programmes for Atong were found in search results. The dictionary below is the nearest thing to a basis for building teaching materials.
\begin{itemize}
\item A Dictionary of Atong, a Tibeto-Burman language of Northeast India and Bangladesh (Seino van Breugel, De Gruyter Mouton, 2021): Atong-English and English-Atong dictionary with semantic and grammatical lexica. Catalogue record: \url{https://oeaw.ac.at/resources/Record/9783110718003}; further record: \url{https://vufind.lib.tsu.ru/Record/libtsu.1013546/Details}
\end{itemize}

\section{Linguistic and Reference Materials}
\begin{itemize}
\item Review of the Atong dictionary: \url{https://www.jbe-platform.com/content/journals/10.1075/ltba.00012.dri}
\item Overview of the language (encyclopedic, secondary source): \url{https://en.wikipedia.org/wiki/Atong_language_(Sino-Tibetan)}
\end{itemize}
Search results also mention a grammar of Atong (Brill, 2014) and a volume of glossed Atong texts (Brill, 2019) by the same author, and a 2009 first-edition Atong-English dictionary (Tura Book Room); no direct URLs for these were seen.

\section{Religious and Community Resources}
None found in search results.

\section{Digital and Audio Resources}
No audio, video or corpus resources were found in search results.

\section{Curriculum Development Context}
Atong has a Latin-script orthography and one substantial modern dictionary, but no classroom materials were found. A curriculum would need to begin with community consultation and with adapting the dictionary and text collections into readers; the neighbouring Garo language may offer models, but that was not investigated here.

\end{document}